\addcontentsline{toc}{chapter}{Abstract}\vspace{-1cm}
%Border
\begin{tikzpicture}[remember picture, overlay]
  \draw[line width = 4pt] ($(current page.north west) + (0.75in,-0.75in)$) rectangle ($(current page.south east) + (-0.75in,0.75in)$);
\end{tikzpicture}
Abstract is essentially a complete overview or the Synopsis of the project in one page, should be precise, concise, and clear. It has to be written in past tense as the work is already completed. If abbreviations or acronyms need to be included in the abstract, the first time it is used, it should be written out in the full form and put the abbreviation in brackets. e.g. Magnetic Resonance Imaging (MRI). The abstract may have about 3 paragraphs covering the outline of the work as below:

\textbf{First para:} Brief introduction about the project which should include background, Areas/type of applications, market size (if available) and hence the importance of the domain. Key developments in the domain, unresolved issues or/and new emerging opportunities in the domain areas and hence the problem or the issue that is proposed to be addressed in the project or the objective. If it is theory or simulation based add all the existing theory/algorithms names and what is it that the project is based on, their essentially advantage or drawbacks in brief and hence the motivation to work on the project.

\textbf{Second Para:} Objectives, scope, Methodology of your project including the key theory, design tools or software tools or / and the experimental tools that will be used, the key specifications including constrains or boundary conditions or approximations or assumptions that have been considered for the work. The sequence in which the work has been carried out.

\textbf{Third para:} Key findings, results or outcome of the work including experimental data of the out come or / and percentage in terms of working efficiency of the developed product or system.

In brief the Abstract should be
\begin{itemize}
	\item Complete — it covers all aspects of the project
	\item Concise — it contains no excess wordiness or unnecessary information.
	\item Clear — it is readable, well organized, and not too jargon-laden.
	\item Cohesive — it flows smoothly between the parts.
\end{itemize}

\pagebreak